\section{Literature Review}
\label{sec:lit}

Our empirical design draws on three strands of the housing-forecasting literature: state-level forecastability, predictor heterogeneity across local markets, and the comparison between econometric baselines and machine-learning methods. We first draw on \citet{rapach2009forecasting}, who forecast \textit{state-level} real housing price growth in the United States from 1995 to 2006 using a wide set of state, regional, and national economic predictors. We adopt the same core macroeconomic series (mortgage rate, CPI, unemployment) retrieved from FRED. Their central result is that forecastability differs sharply between \textit{interior} and \textit{coastal} US states: autoregressive and multi-predictor models perform relatively well for many interior states, but all models perform poorly for a group of primarily coastal states that experienced especially strong price growth, suggesting a ``disconnect'' between prices and fundamentals there. We extend the spirit of this comparison to the \textit{metropolitan} level, selecting one coastal market (San Francisco), one Sun Belt market (Austin), and one interior market (Cleveland) to test whether the interior-versus-coastal forecastability gap holds across structurally different metros and into a much later sample period.

We also draw on \citet{kholodilin2014business}, who evaluate a large set of indicators to forecast housing prices and rents across 71 large German cities. Their central finding is that \textit{no single predictor} dominates across all cities and market segments, with the most useful indicators (national business confidence, consumer confidence, price-to-rent ratios) varying by local context. This motivates our use of \textit{both} listing-side variables (new listings, days-pending, price cuts) \textit{and} macroeconomic indicators, as well as our application of SHAP attribution to identify which variables drive the predictions in each metro.

Finally, our choice of competing model architectures is motivated by \citet{plakandaras2015forecasting}, who compare a machine-learning specification (Support Vector Regression with Ensemble Empirical Mode Decomposition) against Random Walk, Bayesian Autoregressive, and Bayesian VAR baselines for the US real house price index, and find that the machine-learning approach substantially outperforms the econometric baselines out of sample. We adopt the same design principle, comparing a linear baseline (OLS), a univariate time-series baseline (ARIMA), a tree-based ensemble (XGBoost), and a small neural network (MLP) on the same test set, but at the metropolitan rather than national level. Our contribution sits at the intersection of these three strands: we implement metro-level housing nowcasting on three structurally different US markets, compare classical econometric and machine-learning models on a common test set, and use SHAP attribution to surface the variables that drive each model's predictions city by city.
